\section{Power sums and trace powers}

\begin{lemma}[Uniqueness of nonzero weights from tail power sums]
    \label{thm:asym_tail_weights}
    \lean{Multiset.eq_of_notMem_zero_of_forall_sum_map_pow_eq,
        Multiset.eq_zero_of_notMem_zero_of_forall_sum_map_pow_eq_zero}
    \leanok
    Let $\Lambda$ and $\Lambda'$ be finite multisets in $\C\setminus\{0\}$.
    If their power sums agree for all sufficiently large integers $N$, then
    $\Lambda=\Lambda'$, including multiplicities. In particular, if the
    power sums of $\Lambda$ vanish for all sufficiently large $N$, then
    $\Lambda$ is empty.
\end{lemma}

\begin{proof}
    \leanok
    List the distinct entries of the two multisets as
    $z_0,\ldots,z_{r-1}$, and let $b_j$ be the difference of their
    multiplicities. For a threshold $N_0$ and each $0\leq i<r$,
    \begin{align}
        \sum_{j=0}^{r-1}(b_jz_j^{N_0})z_j^i & = 0.
        \notag
    \end{align}
    The Vandermonde matrix is invertible, so $b_jz_j^{N_0}=0$.
    Since $z_j\neq0$, all multiplicities agree. Comparison with the empty
    multiset proves the vanishing assertion.
\end{proof}

\begin{lemma}[Power sums of modulus one]
    \label{thm:mpoalg_unimodular_power_sum}
    \lean{Complex.exists_norm_eq_one_sum_pow_eq_pow}
    \leanok
    Let $\mu_1,\ldots,\mu_n$ be nonzero complex numbers with
    $\left|\sum_k\mu_k^L\right|=1$ for every $L\geq1$. Then $n=1$ and
    $|\mu_1|=1$.
\end{lemma}

\begin{proof}
    \leanok
    \uses{thm:asym_tail_weights}
    For $L\geq1$,
    $\sum_{k,l}(\mu_k\bar\mu_l)^L=\left|\sum_k\mu_k^L\right|^2=1$, so the
    $n^2$ nonzero numbers $\mu_k\bar\mu_l$ have the power sums of the single
    number $1$ at every positive exponent. By uniqueness of power sums on a
    tail they coincide with $\{1\}$ as a multiset, so $n^2=1$. The case $L=1$
    gives $|\mu_1|=1$.
\end{proof}

\begin{lemma}[Nilpotency from vanishing trace powers]
    \label{thm:asym_nilpotent_trace_powers}
    \lean{LinearMap.isNilpotent_of_forall_trace_pow_eq_zero}
    \leanok
    An endomorphism of a finite-dimensional complex vector space, all of
    whose positive powers have vanishing trace, is nilpotent.
\end{lemma}

\begin{proof}
    \leanok
    The characteristic polynomial is determined by the power sums of its
    roots (Newton's identities), so vanishing trace powers force it to
    equal $X^{\dim}$; the Cayley--Hamilton theorem then gives nilpotency.
\end{proof}

\begin{lemma}[Constant trace powers are integral]
    \label{thm:wordflags_constant_trace_powers}
    \lean{Matrix.trace_aeval_eq_of_forall_trace_pow_eq,
        Matrix.exists_nat_eq_of_forall_trace_pow_eq}
    \leanok
    Let $N$ be a complex square matrix with $\tr(N^k)=c$ for every integer
    $k>0$. Then every eigenvalue of $N$ is $0$ or $1$, and $c$ is the number
    of eigenvalues equal to $1$, counted with multiplicity; in particular $c$
    is a nonnegative integer.
\end{lemma}

\begin{proof}
    \leanok
    \uses{thm:asym_nilpotent_trace_powers}
    For a polynomial $p$, the trace of $p(N)$ is $p(0)$ times the size of
    $N$ plus $c(p(1)-p(0))$. For $p=(X^2-X)^k$ with $k>0$ both values
    vanish, so every positive power of $N^2-N$ has trace zero and its
    characteristic polynomial is a power of $X$. By spectral mapping every
    eigenvalue $\lambda$ of $N$ satisfies $\lambda^2=\lambda$, and
    $c=\tr N$ is the sum of the eigenvalues.
\end{proof}

\begin{lemma}[Semisimple rank-one Jordan--Chevalley reduction]
    \label{thm:mpu_semisimple_scaled_trace_rank_one}
    \lean{Module.End.IsNilpotent.trace_add_pow_eq_trace_pow_of_commute,
        Module.End.finrank_range_eq_one_of_isSemisimple_of_scaled_trace_pow_eq_one,
        Module.End.IsNilpotent.mul_eq_zero_and_mul_eq_zero_of_commute_of_finrank_range_eq_one,
        Module.End.IsNilpotent.exists_add_pow_eq_right_of_mul_eq_zero,
        Module.End.IsNilpotent.products_eq_zero_and_eventual_add_pow_eq_of_scaled_trace}
    \leanok
    Let $V$ be a finite-dimensional complex vector space and let $S$ be a
    semisimple endomorphism of $V$. Suppose that $c\in\C$ is nonzero and
    \begin{align}
        \tr\bigl((cS)^k\bigr) & =1
        \notag
    \end{align}
    for every integer $k>1$. Then $\dim\operatorname{range}(S)=1$.
    If $N$ is nilpotent and commutes with $S$, then, for every positive
    integer $k$,
    \begin{align}
        \tr\bigl((N+S)^k\bigr) & =\tr(S^k), \notag \\
        NS                     & =SN=0. \notag
    \end{align}
    There is an integer $J>0$ such that $(N+S)^k=S^k$ for every $k\geq J$.
    In \cite[Proposition~20, lines~3839--3866]{Molnar2018SemiInjective},
    this identity provides the generic Jordan--Chevalley step.
\end{lemma}

\begin{proof}
    \leanok
    In any basis of $V$, the matrix $A$ of $cS$ has trace one at every power
    above one, so its characteristic polynomial is $X^{\dim V-1}(X-1)$.
    Since $cS$ is semisimple, its generalized zero-eigenspace is its kernel,
    whose dimension is the root multiplicity of zero, namely $\dim V-1$.
    Rank--nullity and $c\ne0$ give $\dim\operatorname{range}(S)=1$.

    For $k>0$, commutation gives the binomial expansion
    $(N+S)^k=\sum_{m=0}^k \binom{k}{m}N^mS^{k-m}$. Every summand with $m>0$
    is nilpotent and has trace zero, so only the term $m=0$ remains.

    Commutation makes $\operatorname{range}(S)$ invariant under $N$. The
    restriction of $N$ to this one-dimensional space is both nilpotent and a
    scalar, hence zero. Therefore $NS=0$, and commutation also gives $SN=0$.
    By induction $(N+S)^k=N^k+S^k$ for $k>0$, and taking $J$ at least the
    nilpotency index of $N$ removes the first summand.
\end{proof}
